% RESEARCH_DRAFT_SHA256: 355805254aa55b8a1cde92c4b8fbcad783231d7a0d40bbbd09ab10cadc0644ff
% Layout and source-reference wrappers expanded for the web renderer.
\documentclass[11pt]{article}
\usepackage[margin=25mm]{geometry}
\usepackage{amsthm}
\usepackage{fontspec,unicode-math}
\setmainfont{Latin Modern Roman}
\setmathfont{Latin Modern Math}
\usepackage{amsmath,amssymb,mathtools,xurl,hyperref}
\hypersetup{colorlinks=true,urlcolor=blue,linkcolor=blue}
\setcounter{secnumdepth}{0}
\setlength{\parindent}{0pt}
\setlength{\parskip}{5pt}
\setlength{\emergencystretch}{3em}
\newcommand{\sref}[2]{\hyperlink{source:#1:#2}{[S#2]}}
\newcommand{\cataloguescope}[1]{\paragraph{Source record.}#1}
\newcommand{\floor}[1]{\left\lfloor#1\right\rfloor}
\newtheorem{proposition}{Proposition}
\newtheorem{lemma}[proposition]{Lemma}
\begin{document}
\paragraph{Research manuscript.} OpenAI \texttt{gpt-6-astra}, reasoning \texttt{ultra}; 11 October 2026.
\paragraph{Preserved source body.} The complete catalogue mathematical text and sources follow with layout and source links expanded, including source qualifications and its historical preparation notes. New research begins at the section headed Research attempt.

% UnsolvedMath ID 30006670; source catalogue code OWR-14299916-008
\section{Duplicate Human-Alignment Question}\label{30006670}
{\small UnsolvedMath ID 30006670 · Computer Science}
\subsection{Definitions and mathematical statement}

\paragraph{Shared notation.}Standard number systems are $\mathbb N,\mathbb Z,\mathbb Q,\mathbb R,\mathbb C$. Specialized notation and index ranges follow the cited source and the definitions below.


A representation maps observations into features used for subsequent decisions. Alignment with human behavior can be evaluated only after specifying human judgments, a comparison measure, and the tasks being compared.
\paragraph{Record scope and qualifications.}
The catalogue title is retained as metadata. Named certificates, contracts, bounds or frameworks in that title are not asserted as results of the original source extract.
The supplied registry flags this extraction as unrecoverable. It is displayed as a literal source fragment; the definitions provide context rather than a replacement formal statement.

Source-record limitation: Source flags the statement as unrecoverable: This extraction repeats the same human-alignment question retained in OWR-14299916-007 and appends an author name and page footer; it is not a distinct problem. No separate clean mathematical statement is supplied; this source ID preserves its provenance.

Related catalogue records: \href{https://mehmetmars7.github.io/Erdosproblems-llm-hunter/problem.html?type=open_problems&id=30006669}{30006669}. Separate source IDs are preserved for their provenance.
The source text below is rendered literally where necessary; damaged formulas and composite extracts are not silently repaired into a new mathematical claim.
The supplied catalogue marks this record as unrecoverable. Its extract and terminology are documented, but a complete original problem has not been reconstructed.
\paragraph{Statement recorded in the supplied source.}
What is missing for modern learning systems to exhibit more human-like behavior? How can the behavior of AI systems be more closely aligned with that of human users? Klaus-Robert Müller Modern and Emerging Phenomena in Machine Learning 713 \href{https://doi.org/10.4171/owr/2026/11}{[S2]}
\paragraph{Source status.}
The supplied catalogue labels this source record partially solved; that metadata was checked against the supplied snapshot on 2026-10-11. The label is preserved for provenance and is not an independent verification of present open status or mathematical correctness. The catalogue literature assessment is dated 2026-08-22.
\subsection{Short English statement}
The supplied fragment concerns duplicate Human-Alignment Question. Its complete mathematical target is not recovered.
\subsection{Sources}
\begin{itemize}
\item[S1] Ulam AI and the UnsolvedMath Contributors, UnsolvedMath: A Curated Collection of Open Mathematics Problems, record 30006670 (OWR-14299916-008). Catalogue curation and metadata: CC BY 4.0. The source extract is retained; definitions, English summaries and qualifications are editorial additions.
\url{https://huggingface.co/datasets/ulamai/UnsolvedMath}.
\item[S2] Modern and Emerging Phenomena in Machine Learning (2026), Oberwolfach Reports, DOI 10.4171/owr/2026/11, pp. 683--736
\url{https://doi.org/10.4171/owr/2026/11}.
\item[S3] Muttenthaler et al., Aligning Machine and Human Visual Representations
\url{https://arxiv.org/abs/2409.06509}.
\item Source-specific Oberwolfach Reports licences, including CC BY-SA 4.0 where applicable, are retained; see THIRD\_PARTY\_NOTICES.md.
\end{itemize}
\paragraph{Preparation and review.}
The original source and UnsolvedMath contributors retain attribution for the source extract. Codex (GPT-6; exact serving revision not exposed) assisted with definitions, English summaries, scope notes and formatting. No new human mathematical authorship or independently verified solution is asserted. Permanent source ID 30006670 and code OWR-14299916-008 checked against the complete supplied registry and current repository. Distinct upstream records, including repeated questions and broad or damaged extracts, are retained for complete coverage of the source CS category. A repeated record is not a new independent conjecture. Complete mathematical-equivalence and current-literature review has not been performed. Related canonical identities: 30006669. Preparation and final review follow contribution prompts C and D, with source gaps disclosed. Structural and rendering checks do not establish a new theorem or open status.
\label{end:30006670}

\section{Research attempt}
\subsection{Literature, source context and scope}
The full catalogue body, including qualifications and historical source status, is preserved above. The original workshop report was inspected in the locally cached publisher PDF text. This attempt distinguishes the source question from any explicitly restricted model introduced below. Repeated records retain their own source IDs and do not become independent conjectures. Elementary results and attributed literature results are not claimed as new original-target progress.

\subsection{A separate uncertainty-identification route for the repeated record}
The repeated source question includes agreement with human uncertainty. We test whether matching all ordinal geometric comparisons of a representation fixes the uncertainty probabilities. This is an explicitly limited representation model; no empirical human behavior is inferred.
Let a finite collection of representations be points $z_i$ in Euclidean space. A triplet comparison records only whether $\|z_i-z_j\|<\|z_i-z_k\|$. Suppose a two-choice model converts squared distances $d_1^2,d_2^2$ into probabilities proportional to $e^{-d_1^2},e^{-d_2^2}$.
\begin{proposition}
All ordinal triplet comparisons are unchanged by a common positive scale factor, but the resulting two-choice uncertainty need not be unchanged and can vary from arbitrarily near $1/2$ to arbitrarily near one.
\end{proposition}
\begin{proof}
Multiplying all points by $c>0$ multiplies every distance by $c$, preserving all strict inequalities and ties. For a query at zero with candidate points zero and one on the real line, scaling gives squared distances zero and $c^2$. The first candidate's probability is
\[
 p(c)=\frac1{1+e^{-c^2}}.
\]
As $c\downarrow0$ this approaches $1/2$, and as $c\to\infty$ it approaches one. The ordinal preference for the zero-distance candidate is the same for every positive $c$.
\end{proof}
Thus ordinal behavioral matching alone leaves a scale or temperature unidentifiable in this probability model. One would need probability-valued judgments or a separately justified scale to claim calibrated uncertainty. This does not say the source's empirical method used only ordinal observations; it rejects the proposed inference without such an additional premise.
\subsection{Verification and remaining gaps}
At $c=1$ and $c=2$, the probabilities are approximately $0.731059$ and $0.982014$, despite identical ordinal geometry. Representation transfer, fine-to-coarse structure and human uncertainty calibration remain untested. This repeated record is not treated as a separate solved conjecture.

\section{References and source attribution}
The primary Oberwolfach workshop report, \url{https://doi.org/10.4171/owr/2026/11}, was inspected on 11 October 2026. Relevant source passages: 2026/11 pp. 712--713, repeated human-alignment and uncertainty question. Catalogue and contributor attribution remain in the original Sources section. Mathematical arguments labeled as this attempt are proved above; cited external results remain attributed to their authors and are not independently reproved or claimed as our discoveries.

\url{https://arxiv.org/abs/2409.06509}. Muttenthaler and coauthors, Aligning Machine and Human Visual Representations across Abstraction Levels. The primary abstract was opened and its task-specific representation-transfer context checked. Its empirical improvements are not claimed as findings of this attempt, which collected no human-judgment data.


\paragraph{Integration review.} Codex (GPT-6; exact serving revision not exposed) reviewed the written arguments and expanded layout and source-reference wrappers. Repository metadata and the production web renderer were checked using Python 3.13 and JavaScript with MathJax 3.2.2. This model-assisted review is not independent human proof verification; no PDF compilation or formal proof checking is asserted.

\section{Completion Estimate}
\noindent\textbf{COMPLETION: 0\%}
No substantiated new progress toward the entire original target is claimed. The worked results and failed route are delimited above; a known result, illustrative model, or auxiliary obstruction is not counted as original-target completion.

\section{Final status and attribution}
\textbf{Attempt status: unresolved.} This describes this research attempt, not a new ruling on the source catalogue's status or on external literature claims.
The mathematics and this self-audit were generated by OpenAI \texttt{gpt-6-astra}, reasoning \texttt{ultra}, on 11 October 2026. Codex (GPT-6; exact serving revision not exposed) orchestrated the batch and contributed formatting/repository review. No human mathematical authorship, independent human proof checking, or formal verification is claimed. The historical preparation paragraphs in the reproduced catalogue body concern that body only.

Contribution prompt A was applied to the complete definition and supplied source record, followed by the adversarial checks of prompt B. Local reproduction files are \texttt{batch011/checks.py}, \texttt{computations.json}, and the per-ID audit JSON. They were not public links at preparation time. Numerical checks are finite and do not prove unrestricted results. No LaTeX engine was found locally; compilation was not performed by this research agent.

\end{document}
